\section{Related Work}
\label{sec:related}

\paragraph{Agentic task benchmarks.}
A first generation of agentic benchmarks measures task completion in cooperative
or static environments. SWE-bench evaluates code-patch generation against real
GitHub issues with executable test oracles \citep{jimenez2024swebench}; WebArena
and GAIA measure web navigation and general assistant competence
\citep{zhou2024webarena,mialon2023gaia}; AgentBench aggregates many such
environments \citep{liu2023agentbench}. These provide \emph{verifiable} success
signals but assume the environment is not an adversary. VEB retains verifiability
of \emph{sub-goals} (a fact is disclosed or it is not; an economic-buyer meeting
is earned or it is not) while making the counterpart strategically
non-cooperative.

\paragraph{Simulated-user and tool-agent dialogue.}
$\tau$-bench pairs a tool-using agent with an LLM-simulated user and grades the
final database state against a policy \citep{yao2024taubench}, and ALFWorld and
related work embed agents in simulated worlds \citep{shridhar2021alfworld}. VEB
shares the simulated-interlocutor design but differs in two ways: the interlocutor
is a \emph{committee} of personas with conflicting incentives and gated private
knowledge, and grading targets the \emph{quality of the persuasion process}, not
only a terminal state.

\paragraph{Negotiation, persuasion, and social simulation.}
Prior work studies LLMs in bargaining and dialogue games---deal-or-no-deal
negotiation \citep{lewis2017dealornodeal}, diplomacy-style strategic play
\citep{fair2022diplomacy}, and persuasion dialogue \citep{wang2019persuasion}.
More recently, NegotiationArena evaluates LLM-vs-LLM bargaining across structured
negotiation games \citep{bianchi2024negotiationarena}, and Sotopia evaluates
social intelligence in open-ended goal-driven interactions between LLM agents
\citep{zhou2024sotopia}. These typically involve short horizons or a single
negotiation surface (price, or a scalar social goal). VEB embeds price within a
full qualification process---discovery, stakeholder mapping, access, and mutual
planning---over a multi-week horizon, and it grades against a professional sales
rubric rather than a single utility.

\paragraph{Sales-domain agent benchmarks.}
Closest in domain is CRMArena \citep{huang2025crmarena}, which evaluates agents
on realistic CRM \emph{operational} tasks (case routing, analytics, policy
compliance) inside a simulated Salesforce org. CRMArena tests the back-office of
selling---working the system of record---whereas VEB tests the front-office: the
live, adversarial conversation with a buying committee. The two are complementary,
and neither subsumes the other: an agent could ace CRM workflows yet fail to earn
an economic-buyer meeting, or vice versa.

\paragraph{LLM-as-judge and its validation.}
Using strong LLMs to grade open-ended outputs is now standard
\citep{zheng2023judge}, with known biases (position, verbosity, self-preference)
and a corresponding need for human calibration. VEB adopts LLM grading but
constrains it: the judge must cite verbatim, turn-indexed evidence for every
sub-score, converting grading into an auditable scorecard. We pair this with an
explicit human-expert validation protocol whose pre-registered inter-rater
agreement study is forthcoming (Section~\ref{sec:validation}), following best
practice for defensible LLM-graded benchmarks.

\paragraph{Preference-based post-training.}
Reinforcement learning from human (or AI) feedback aligns models to preferences
\citep{ouyang2022instructgpt,bai2022constitutional}; DPO learns directly from
preference pairs without an explicit reward model \citep{rafailov2023dpo}; and
process reward models supervise intermediate reasoning steps
\citep{lightman2023letsverify}. VEB is designed to \emph{produce} the data these
methods need: its seed-matched win/loss pairs and turn-level, evidence-cited
scorecards constitute preference pairs and process labels over a verifiable
environment, making VEB usable not only for evaluation but for training
negotiation-capable agents.
